\documentclass{article}

\usepackage{graphicx}
\usepackage{listings}

\begin{document}

\begin{center}


\includegraphics[width=3cm]{iiitb_comet_logo.jpg}

\vspace{0.4cm}

IIIT Bangalore

\vspace{0.3cm}

Name: Srinidhi A

ID: COMETFWC053

\vspace{0.5cm}

Question 36

\end{center}

\section*{Question}

Simplify the Boolean expression and determine the logic function.

\section*{Analysis}

For inputs $Y$ and $Z$ the function is

$F = \overline{YZ}$

This represents the NAND operation.

Truth table:

\begin{tabular}{c c c}
$Y$ & $Z$ & $F$ \\
0 & 0 & 1 \\
0 & 1 & 1 \\
1 & 0 & 1 \\
1 & 1 & 0 \\
\end{tabular}

\section*{Python Code}

\begin{lstlisting}
import numpy as np
import matplotlib.pyplot as plt

Y = np.array([0,0,1,1])
Z = np.array([0,1,0,1])

F = 1 - (Y & Z)

print("Y Z F")

for i in range(4):
    print(Y[i], Z[i], F[i])

x = np.arange(4)

plt.step(x,Y,where='post',label='Y')
plt.step(x,Z,where='post',label='Z')
plt.step(x,F,where='post',label='F')

plt.grid()
plt.legend()

plt.savefig("graph.png")

plt.show()
\end{lstlisting}

\section*{Graph}

\begin{center}

\includegraphics[width=10cm]{graph.png}

\end{center}

\section*{Conclusion}

The output follows the NAND truth table. Therefore the simplified Boolean expression corresponds to a NAND gate.

\end{document}
